% !TEX program = xelatex
\documentclass[a4paper,11pt]{article}
\usepackage{fontspec}
\usepackage{xeCJK}
\setCJKmainfont[AutoFakeBold=2]{Droid Sans Fallback}
\setCJKsansfont[AutoFakeBold=2]{Droid Sans Fallback}
\setCJKmonofont[AutoFakeBold=2]{Droid Sans Fallback}
\usepackage{amsmath}
\usepackage{array}
\usepackage{geometry}
\usepackage{longtable}
\usepackage{xcolor}
\usepackage{booktabs}
\usepackage{enumitem}
\usepackage{fancyhdr}
\usepackage{amssymb}
\usepackage{hyperref}

\geometry{top=2.35cm,bottom=2cm,left=1.8cm,right=1.8cm,headheight=24pt,headsep=6pt,footskip=1cm}
\pagestyle{fancy}
\fancyhf{}
\fancyhead[L]{\small CSP-J 入门级知识清单}
\fancyhead[R]{\small\nouppercase{\leftmark}}
\fancyfoot[C]{\thepage}
\renewcommand{\headrulewidth}{0.4pt}
\renewcommand{\footrulewidth}{0pt}

\newcolumntype{L}[1]{>{\raggedright\arraybackslash}p{#1}}
\newcolumntype{C}[1]{>{\centering\arraybackslash}p{#1}}
\newcolumntype{R}[1]{>{\raggedleft\arraybackslash}p{#1}}
\newcommand{\code}[1]{\texttt{#1}}
\setlist[itemize]{nosep,leftmargin=1.5em}
\setlist[enumerate]{nosep,leftmargin=1.8em}
\setlength{\tabcolsep}{3pt}
\renewcommand{\arraystretch}{1.12}
\setlength{\LTleft}{0pt}
\setlength{\LTright}{0pt}
\setlength{\LTpre}{5pt}
\setlength{\LTpost}{5pt}
\setlength{\parindent}{0pt}
\setlength{\parskip}{2pt}
\emergencystretch=2em
\hbadness=10000
\renewcommand{\contentsname}{目录}
\hypersetup{
  colorlinks=true,
  linkcolor=black,
  urlcolor=blue,
  pdftitle={CSP-J入门级知识清单},
  pdfauthor={scmmm}
}

\begin{document}

\thispagestyle{empty}
\begin{center}
{\LARGE\bfseries CSP-J 入门级知识清单} \\
\vspace{3mm}
{\large 知识范围、复杂度与实现检查} \\
\vspace{3mm}
{\small 示例基线：C++14（\code{-O2 -std=c++14}），具体语言版本、时限和内存以当年比赛公告为准。}
\end{center}

{\small\textbf{使用说明：}逐项确认适用条件、正确实现和复杂度，并为关键结论准备边界数据或反例。}

{\small\textbf{星号说明：}带 * 的内容不属于入门级大纲的知识类别，属于提高级、NOI 级或补充内容；入门级类别下的常见实现、变体和应用不标星。}

\vspace{2mm}
\tableofcontents
\newpage

%==================================================================
\section{入门级知识范围核对}
%==================================================================

本节列出入门级需要掌握的基础知识与对应检查项。

\begingroup
\small
\begin{longtable}{L{2.5cm} L{8.6cm} L{5.5cm}}
\toprule
知识模块 & 必须掌握的完整范围 & 检查要求 \\
\midrule
\endfirsthead
\toprule
知识模块 & 必须掌握的完整范围 & 检查要求 \\
\midrule
\endhead
\endfoot
\bottomrule
\endlastfoot

基础与环境
& CPU、内存、I/O；位、字节与字；Windows/Linux 与文件目录操作；网络和 Internet 基本概念；计算机与 NOI 历史、活动规则；语言、编译、运行、IDE 与 g++。
& $\square$ 把 512 MiB 换成字节和 bit；不用 IDE 说明源码怎样编译、运行和重定向。 \\[2mm]

程序基础
& 标识符、关键字、常量、变量、字符串、表达式；命名与作用；头文件、名字空间；编辑/编译/解释/调试；四类基本类型；输入输出、条件、循环；算术/关系/逻辑/自增/三目/位运算；数学库；模块化设计与流程图。
& $\square$ 区分声明、初始化和赋值；把含循环、无解分支的流程图转成 C++14。 \\[2mm]

数组与函数
& 一维/多维数组；字符数组与 string；函数定义和调用、形参与实参、值/引用传参、作用域与递归。
& $\square$ 计算多维数组字节数；读取含空格行；把复制 vector 的函数改成 \code{const\&}。 \\[2mm]

复合类型
& 结构体、联合体、枚举；指针、指针数组访问、字符指针、结构体指针和引用。
& $\square$ 用结构体和枚举描述搜索状态；判断指针/引用在函数返回和 vector 扩容后是否有效。 \\[2mm]

文件与 STL
& 文本/二进制文件、重定向和文件读写；min/max/swap/sort；stack/queue/list/vector。
& $\square$ 说明文本读数与读取二进制表示的差别；解释 list 无随机访问且不能直接套 \code{std::sort}。 \\[2mm]

线性结构
& 单链表、双向链表、循环链表、栈和队列；链表 O(1) 插删的前提是已经持有正确节点。
& $\square$ 画出删除双向链表中间节点前后的四个指针变化。 \\[2mm]

简单树
& 树与二叉树的定义、性质、表示和存储；二叉树前序、中序、后序遍历。
& $\square$ 对含空孩子的树写三种遍历，并解释一种遍历为何通常不能唯一还原树。 \\[2mm]

特殊树
& 完全二叉树及数组表示；哈夫曼树的构造与编码；二叉搜索树的定义和构造。
& $\square$ 区分完全/满/BST；用递增插入说明普通 BST 可退化为 O($n$)。 \\[2mm]

简单图
& 图的点、边、方向、度、路径、连通等概念；邻接矩阵和邻接表。
& $\square$ 比较 O($V^2$) 与 O($V+E$) 空间，并正确处理无向边的两条弧。 \\[2mm]

算法与策略
& 算法概念；自然语言、流程图、伪代码描述；枚举、模拟、贪心、递推、递归、二分、倍增；前缀和与差分。
& $\square$ 为贪心写证明、为二分写单调谓词、为枚举证明不重不漏。 \\[2mm]

高精度
& 高精度加法、减法、乘法；高精度整数除单精度整数的商和余数。
& $\square$ 用 \code{1000-1}、\code{999*999}、\code{1000/7} 检查借位、进位和余数。 \\[2mm]

排序
& 排序概念；冒泡、选择、插入、计数排序；稳定性、时间和额外空间。
& $\square$ 在重复键数据上判断四种排序的稳定性；说明计数排序依赖值域。 \\[2mm]

搜索与图遍历
& DFS、BFS；图的深度/广度优先遍历；Flood Fill。
& $\square$ 给带权图反例说明普通 BFS 不能求一般最短路；在含环图中保证状态只访问一次。 \\[2mm]

动态规划
& DP 基本思想；简单一维、简单背包、简单区间 DP。
& $\square$ 分别写状态、初值、转移顺序和答案位置；不可达状态不能默认设 0。 \\[2mm]

数与数论
& 数集和四则运算；二/八/十/十六进制；初中代数/几何；整除、因数、倍数、指数、质合数、取整、模、唯一分解、欧几里得、埃氏筛和线性筛。
& $\square$ 测试 0/1/2 与平方数；区分 C++ 负数除法和数学 floor。 \\[2mm]

组合与其他
& 集合；加法/乘法原理；排列、组合、杨辉三角；ASCII 码。
& $\square$ 用小枚举检查计数是否重复；不用魔法数字完成数字字符转换和大小写判断。 \\
\end{longtable}
\endgroup

\newpage

%==================================================================
\section{C++ 语言与实现基础}
%==================================================================

\begin{longtable}{L{2.4cm} L{7.5cm} L{6.8cm}}
\toprule
知识点 & 适用条件 & 检查要求 \\
\midrule
\endfirsthead
\toprule
知识点 & 适用条件 & 检查要求 \\
\midrule
\endhead
\endfoot
\bottomrule
\endlastfoot

整数类型范围与溢出
& 用于可能产生大整数的运算。类型选择须覆盖输入、答案和中间结果；扩大目标变量类型不能修复已经发生的窄类型溢出。
& 实现受限乘方：输入 $a,b$（$a,b\le10^9$），若 $a^b>10^9$ 输出 $-1$；扩型和溢出检查须在乘法之前完成。 \\[2mm]

整数除法与取模
& 整除、向上取整、余数状态和二进制拆分。C++14 的有符号整数除法向零取整；负数取模的符号也要按语言规则判断。
& 给定 $a=-7,b=3$，输出 C++14 下 \code{a/b} 和 \code{a\%b}，解释结果差异。 \\[2mm]

浮点误差与精确比较
& 百分比、$\sqrt{\Delta}$ 和分数比较优先使用整数。\code{double} 无法精确表示 $0.1$。
& 实现整数二分开方 \code{isqrt(n)}，不使用 \code{sqrt}；用 $n=10^{18}$ 验证。 \\[2mm]

数组、字符、string
& 读含空格的句子、解析 IP、统计字符频次。\code{cin >> s} 会跳过空格。
& 读一整行（可能含空格和数字），统计字母和数字各出现几次。 \\[2mm]

函数、引用、指针
& 传递大型对象（如 \code{vector}、邻接表）时避免不必要复制；小型标量按值传递未必更慢。
& 将 \code{void f(vector<int> v)} 改为 \code{void f(const vector<int>\& v)}，解释快的原因。 \\[2mm]

递归与调用栈
& 递归适用于树遍历、分治排序和表达式求值；深度较大时可能耗尽调用栈。栈帧大小和栈上限取决于编译器、平台和局部变量，须按目标环境核对。
& 分别实现递归 DFS 和显式栈 DFS，并根据最大深度判断是否需要迭代实现。 \\[2mm]

输入输出与文件 IO
& 输入输出量大时可用 \code{std::ios::sync\_with\_stdio(false)} 与 \code{std::cin.tie(nullptr)}；是否超时取决于数据量、输出量和时限，不存在固定的 $n$ 阈值。
& 用 \code{freopen} 重定向的例子。解释 \code{cin >> } 后接 \code{getline} 读到空串的原因。 \\[2mm]

vector、queue、*deque
& 邻接表、BFS 队列、模拟队列。分清各容器适用场景。
& 手写循环队列（\code{push/pop/front}，用数组），说出 \code{std::queue} 底层是 \code{deque}。 \\[2mm]

*STL map、set、heap
& 适用于值域较大、需要有序遍历或需要反复取得最小距离的情况。\code{priority\_queue} 默认按最大元素优先。
& 实现最小优先队列，并据此完成 Dijkstra。 \\[2mm]

运算符优先级、结合性与短路
& 写表达式解析或混合多个运算符时，不能靠猜优先级。位运算 \code{\& | \^{}} 优先级低于 \code{== !=}；逻辑 \code{\&\& ||} 有短路语义；赋值 \code{=} 优先级极低。括号是最便宜的正确性保障。
& 给出 \code{int a=2,b=3,c=4;}，不加括号解释 \code{a\&b==2} 和 \code{a|b\&c} 的实际求值顺序，并说明为何 \code{if(x\&1==0)} 不判断 x 最低位是否为 $0$。 \\[2mm]

C++14 版本与编译选项
& 本文以 \code{-std=c++14} 为基线；\code{std::gcd}、结构化绑定、\code{if constexpr} 属于 C++17，不能依赖编译器扩展。
& 用 \code{g++ -std=c++14 -O2} 编译最近写的一份代码；静态链接、编译器和内存限制以比赛环境为准。 \\
\end{longtable}

\newpage
%==================================================================
\section{基本问题求解与复杂度}
%==================================================================

\begin{longtable}{L{2.4cm} L{7.5cm} L{6.8cm}}
\toprule
知识点 & 适用条件 & 检查要求 \\
\midrule
\endfirsthead
\toprule
知识点 & 适用条件 & 检查要求 \\
\midrule
\endhead
\endfoot
\bottomrule
\endlastfoot

模拟
& 适用于状态转移和操作顺序已由题面明确给出的情形。实现时须完整保留条件和执行顺序。
& 将题面描述拆成有先后约束的单步操作，并标明每步修改的状态。 \\[2mm]

枚举
& 数据规模小的时候枚举所有可能性；剪枝只能降低实际运行量，不能改变最坏复杂度。
& 写 $N$ 个数的子集枚举（位运算版和递归 DFS 版各一）。 \\[2mm]

递归、*分治、显式栈
& 适用于分治排序、树遍历和括号解析；最大深度较大时应采用迭代实现。
& 分别实现递归版和迭代版归并排序，说明递归深度与栈空间差异。 \\[2mm]

贪心
& 适用于局部选择可通过交换论证、支配关系或不变量推出全局最优的情形。
& 对一个未经证明的局部最优策略构造反例，再补充成立条件或改用其他算法。 \\[2mm]

*贡献计算
& 把总答案拆成每个元素、位置、边或区间的独立贡献，常能把两重枚举改成一次统计、前缀和或频次维护。
& 求所有子数组元素和时，推导 $a_i$ 被多少个子数组包含；再用小规模暴力对拍贡献式。 \\[2mm]

前缀和、差分
& 前缀和适合频繁区间聚合查询；差分适合区间加减，最后扫描一次还原。两者的 O(1) 是针对特定操作，不是所有区间操作。
& 实现数组区间加 $1$、最后输出所有位置的值，要求 O($N+Q$)。 \\[2mm]

二分查找与二分答案
& 有序数组中找边界，或答案范围大（如 $10^9$）且判定具有单调性；先区分找任意值、第一处真和最后一处真。
& 分别写 \code{lower\_bound}、第一处可行和最后处可行模板，处理空区间、重复值和中点溢出。 \\[2mm]

*三分查找
& 目标函数在区间内单峰或单谷，且能比较两个位置的函数值时使用；不满足单峰性不能套模板。
& 在整数区间上写三分并处理平台、端点和最终枚举；说明实数三分的精度停止条件。 \\[2mm]

双指针、滑动窗口、*单调队列
& 时间窗（24 小时内乘客）、价格窗（45 分钟票）、坐标窗口：窗口单向滑动。
& 证明"每个元素最多出入队一次"。举看似两重循环、实为 O($N$) 的例子。 \\[2mm]

*离散化
& 值域为 $10^9$ 而有效值只有 $10^5$ 个时，按原值域分配数组会超出常见内存限制。
& 整数数组排序去重后用 \code{lower\_bound} 做压缩映射，写出完整代码。 \\[2mm]

倍增
& 树上跳 $k$ 级祖先、数组跳区间。把 $k$ 按二进制拆位。
& 给出 $up[j][v]$ 的两种含义：从 $v$ 跳 $2^j$ 步到哪，以及转移方程。 \\[2mm]

*复杂度分析
& 统计所有循环和数据结构操作的总执行次数。队列和栈常按每个元素的总进出次数进行摊销分析。
& 给一段同时含 \code{while} 弹队列和 \code{for} 扫描的代码，判断 O($N$) 还是 O($N^2$)。 \\[2mm]

*测试点与部分分
& 若题面提供测试点或分组限制，小规模、全相等、$k=0$ 等条件可能对应部分分；没有测试点表时也要根据约束自行分层。
& 根据题目约束与测试点，分别标注可用直接算法和需要优化的子任务。 \\[2mm]

*边界数据
& 补充样例未覆盖的最小值、最大值、全相等、无解和等号边界。
& 为一段排序代码造 10 组数据，覆盖升序、降序、全相等、含负数、单元素。 \\[2mm]

*正确性说明
& 为核心状态和循环写出不变量，核对初始化、更新顺序和结束条件。
& 说明核心循环执行前后关键变量的含义，并据此检查边界。 \\
\end{longtable}

\newpage
%==================================================================
\section{字符串}
%==================================================================

\begin{longtable}{L{2.4cm} L{7.5cm} L{6.8cm}}
\toprule
知识点 & 适用条件 & 检查要求 \\
\midrule
\endfirsthead
\toprule
知识点 & 适用条件 & 检查要求 \\
\midrule
\endhead
\endfoot
\bottomrule
\endlastfoot

字符串解析与格式转化
& 进制转换（读入 $k$ 进制字符串转整数、或整数转 $k$ 进制输出）、高精度读入（逐字符转数字数组）、IPv4/端口格式校验（\code{a.b.c.d:e} 分五段，禁前导零、检查值域）。核心是逐字符处理、手工分段、用整数数组承接。
& 写 $k$ 进制转十进制：逐位累加 \code{ans = ans * k + digit}。写高精度读入：把长度 $\le10^6$ 的数字串转成 \code{vector<int>}。写 IPv4 解析：严格拒绝空段、前导零（如 \code{01}）、越界（$0\sim255$ 和 $0\sim65535$）和多余字符。 \\[2mm]

字符串直接扫描
& 当 $n,q\le1000$ 且 $nq\le5\times10^8$ 时，可优先检查逐对比较子串或逐查询扫描全表；是否覆盖全部测试点由完整约束决定。
& 实现后缀查询：每次扫描全部书号，用 \code{code \% 10\^{}len == need} 判断后缀，并写出 O($nq$) 时间和 O(1) 额外空间。 \\[2mm]

回文判断、子串与字符频次
& 回文可用 O($N$) 双指针或构造候选年份；子串计数可用滑动窗口或枚举起点；字符频次可用固定大小数组统计。
& 写回文日期题：枚举 4 位年份按回文规则构造 8 位日期，验证月日合法性后判断是否在输入区间内。再写标题统计题：逐字符分类计数字母和数字，注意空格和换行。说明"枚举年份构造"和"逐日推进"的复杂度差异。 \\[2mm]

*Trie（前缀树）
& 大量字符串、前缀查询；即使不直接裸考，也可能作为字典、前缀统计或状态结构的组件。
& 用数组写 Trie，插入 \code{cat/car/dog}，查前缀 \code{ca} 的词数。 \\[2mm]

*字符串哈希
& 快速比较子串、判断树镜像。预处理 O($L$) 后每次查询 O(1)，但哈希存在碰撞风险。
& 约定区间为半开区间 $[l,r)$，推导 $\text{hash}(l,r)=H[r]-H[l]\times\text{base}^{\,r-l}$；说明模数、底数和双哈希如何降低风险。 \\
\end{longtable}

\newpage
%==================================================================
\section{数据结构、搜索与图论}
%==================================================================

\begin{longtable}{L{2.4cm} L{7.5cm} L{6.8cm}}
\toprule
知识点 & 适用条件 & 检查要求 \\
\midrule
\endfirsthead
\toprule
知识点 & 适用条件 & 检查要求 \\
\midrule
\endhead
\endfoot
\bottomrule
\endlastfoot

栈、队列、*双端队列
& BFS（队列）、括号匹配（栈）、滑动窗口最值（单调队列）。元素最多进出一次是摊销核心。
& 手写循环队列（数组实现，不用 STL），解释判空判满为什么要留一个空位。 \\[2mm]

*堆与优先队列
& Dijkstra 取最小距离、事件模拟最早发生。默认大根堆，取最小值得改比较器。
& 写 Dijkstra：为什么同一点可能入堆多次？旧的必须跳过？ \\[2mm]

*集合、映射、哈希
& 值域小用数组；需要有序遍历或最坏复杂度保证时用 \code{map}；能接受期望 O(1) 且注意退化风险时再用 \code{unordered\_map}。
& 构造一组让 \code{unordered\_map} 退化成 O($N$) 的输入。 \\[2mm]

*多重计数与值映射
& 重复值按频次、类别或出现位置统计；值域小用频次数组，值域大用 \code{map}、哈希或离散化，不能把"不同值个数"和"元素总数"混为一谈。
& 对 $[3,1,3,2,3]$ 同时输出每个值的频次、出现位置和不同值个数，比较数组、\code{map} 与离散化三种写法。 \\[2mm]

*并查集
& 只问"两个点是否连通"，不问路径。
& 手写 \code{find}（路径压缩）和 \code{merge}（按大小合并），说明为什么路径压缩不破坏正确性。 \\[2mm]

树与二叉树
& 用于树的遍历、子树大小和镜像判断。须统一空子树表示；深度较大时改为迭代实现。
& 用数组存二叉树，写先序和后序遍历，禁止递归。 \\[2mm]

图存储与遍历
& 稀疏图或点数较大时用邻接表；很小的稠密图才考虑邻接矩阵。有向图、无向图建图不同。
& 给一张 $N$ 点 $M$ 边无向图，用邻接表，BFS 求 $1$ 到各点最短步数。 \\[2mm]

链表与*有序块删除
& 反复取走序列元素、合并相邻同类块。链表在已经定位节点时删除可为 O(1)，查找节点和维护块边界仍可能成为瓶颈。
& 模拟 \code{0011100}：每轮取每个块最左元素，写出每轮输出。 \\[2mm]

*树状数组
& 单点修改、前缀和查询。比线段树短。稳定排序里的排名也能算。
& 手写 \code{add} 和 \code{sum}，说明 lowbit。在 $[3,1,3]$ 上模拟全部更新。 \\[2mm]

DFS、BFS
& 可达性、无权最短路、连通块。标记要在入队时做，而非出队时。
& 在有环图上逐层写出 BFS 入队顺序，解释为何标记时机重要。 \\[2mm]

*Dijkstra 与 0-1 BFS
& 边权非负时可用 Dijkstra；边权仅为 $0,1$ 时可用 deque 做 0-1 BFS。一般非负权不应误写成普通 BFS。
& 用天平类比松弛：为什么非负权下当前弹出的最小有效距离可以定型？旧堆项为什么必须跳过？ \\[2mm]

*状态图（奇偶/资源维）
& "恰好 $L$ 步"、"最多 $k$ 次魔法"、"必须模 $k$ 时刻"：加维度进状态。
& 无向图 $s$ 到 $t$ 恰好 $L$ 步不能只看最短路；还要考虑连通分量、二分图奇偶性和是否存在可往返的边或奇环。构造二分图和含奇环的图，分别判断恰好 $L$ 步可达性，说明为什么"来回 $2$ 步"只是常用的扩展手段，不是完整判定。 \\
\end{longtable}

\newpage
%==================================================================
\section{动态规划}
%==================================================================

\begin{longtable}{L{2.4cm} L{7.5cm} L{6.8cm}}
\toprule
知识点 & 适用条件 & 检查要求 \\
\midrule
\endfirsthead
\toprule
知识点 & 适用条件 & 检查要求 \\
\midrule
\endhead
\endfoot
\bottomrule
\endlastfoot

状态定义、转移、初始化
& 写 DP 先定状态：$dp[i]$ 是"以 $i$ 结尾"还是"前 $i$ 个中选"。初始化不能默认 $0$；"恰好"通常用 $-\infty$ 标不可达，"至多"或空集合法可用 $0$。
& 写"恰好 $n$ 根拼数字"的 DP，解释不可达为何不能初始化为 $0$。再写"网格取数"：允许上下右移动、格权可为负，全负网格为何 DP 初值不能是 $0$ 而要用 $-\infty$。 \\[2mm]

DP 识别：从暴力到 DP
& 题面要求"最多/最少/方案数"、数据 $n\le500$ 但 $2^n$ 枚举不可能时，优先考虑 DP。核心信号：存在重叠子问题（同一状态被多次计算）、存在最优子结构（大问题的最优解包含子问题的最优解）。
& 给摆渡车问题：$n$ 个同学在不同时刻到达，发车间隔固定为 $m$。先写按到达时间分组的 O($2^n$) 暴力枚举发车方案，再指出"前 $i$ 个人的最小等待和"只依赖前面状态，写出 $dp[i]$ 转移。对比暴力与 DP 的复杂度差距。 \\[2mm]

背包 DP
& 选若干物品使和满足某条件。对正权容量的一维 DP，0/1 背包通常从大到小，完全背包通常从小到大；计数、负权或其他转移要重新证明顺序。
& 写 0/1 背包，解释容量循环方向，再改成完全背包。 \\[2mm]

*网格 DP / 带预算 DP
& 网格上单向走、带预算 $k$ 选点。$1000\times1000$ 时不能 O($N^3$)。预算维通常来自题面的 $k\le100$ 这类小参数信号。
& $3\times4$ 网格手动跑 DP，只准右/下，取最大和。再写出状态为 $dp[i][j][used]$ 的三维带预算版。 \\[2mm]

线性 DP 与区间 DP 入门
& 线性 DP（$dp[i]$ 只依赖前若干项）是 CSP-J 最常考的 DP 形态。区间 DP（$dp[l][r]$ 依赖子区间）在回文、括号、取数等场景出现。两者都可用 O($N^2$) 拿部分分，再优化。
& 写最长上升子序列 O($N^2$) 版。再写"网格取数"的按列 DP：每列从上到下、从下到上各扫一次，解释为何不能在同一列内来回走（违反不能重复经过的约束）。 \\[2mm]

*单调队列优化 DP
& 转移从滑动窗口取 max/min 时，直接枚举多 O($W$)，单调队列压到 O($N$)。
& 数组 $[5,1,4,4,2]$，窗口 $3$，逐步写出队列变化。 \\[2mm]

*方案恢复与滚动数组
& 只问最优值时可滚动；要恢复方案时通常保存父指针或决策，不能把必要信息一起覆盖。也可用重算等方法换空间。
& 在背包 DP 上加 \code{choice} 数组，输出最优值的同时输出选了哪些物品。 \\[2mm]

*DP 维度选择与部分分信号
& 从题面约束倒推 DP 维数：$n\le500$ 且 $k\le100$ → 二维 $dp[i][k]$（带预算的点列选择）；$m\le100$ 且 $n\le500$ → $dp[i]$ 一维 + 余数优化（固定间隔发车）；网格 $1000\times1000$ → 按列滚动（网格取数）；值域小（$a_i\le5000$）→ 以值域为容量做背包 DP（选木棍拼多边形）。
& 给一道新题的数据范围，能写出至少两种可能的 DP 状态设计，并估算各自的时间/空间复杂度，选出更可行的方案。 \\
\end{longtable}

\newpage
%==================================================================
\section{数学}
%==================================================================

\begingroup
\small
\renewcommand{\arraystretch}{1.04}
\begin{longtable}{L{2.4cm} L{7.5cm} L{6.8cm}}
\toprule
知识点 & 适用条件 & 检查要求 \\
\midrule
\endfirsthead
\toprule
知识点 & 适用条件 & 检查要求 \\
\midrule
\endhead
\endfoot
\bottomrule
\endlastfoot

基本运算与边界
& max/min、绝对值。\code{abs(INT\_MIN)} 的结果无法由 \code{int} 表示，取负或调用对应重载可能触发未定义行为。
& 不用分支和库函数写整数绝对值。说出 \code{abs(INT\_MIN)} 的后果。 \\[2mm]

位运算
& 判奇偶、提取二进制位、异或（$a\oplus a=0$）、lowbit、枚举子集（\code{for(int s=mask;s;s=(s-1)\&mask)}）。
& 写代码：对正整数只用位运算判 $2$ 的幂，不准用循环；用位运算枚举 $\{a,b,c\}$ 的所有非空子集。 \\[2mm]

质数判定与埃氏筛
& 单个数判质数可试除到 $\sqrt{n}$；需要批量筛一段连续范围的质数时用埃氏筛，O($N\log\log N$)。
& 手写埃氏筛求 $[1,100]$ 的质数表，说明为什么内层从 $i^2$ 开始、为什么只用 $\le\sqrt{N}$ 的质数去筛。 \\[2mm]

线性筛（欧拉筛）与最小质因子
& 要求 O($N$) 或需要每个数的最小质因子时用线性筛。核心：每个合数只被其最小质因子筛掉一次。可顺势维护 \code{lp[x]} 用于 O($\log x$) 因数分解或积性函数递推。
& 写线性筛同时维护 \code{prime[]} 和 \code{lp[x]}，在纸上模拟 $\{2\ldots30\}$：指出 12 被 2 筛、15 被 3 筛、25 被 5 筛，说明判断 \code{i\%prime[j]==0} 后 \code{break} 的原因。 \\[2mm]

gcd 与 lcm
& 约分、分数化简。\code{std::gcd} 是 C++17 才有，C++14 自己写；计算 lcm 时先除后乘以降低溢出风险。
& 手算 $\gcd(12345,67890)$，写每一步余数。 \\[2mm]

*模运算与快速幂
& 答案模 $10^9+7$ 或 $998244353$；指数很大时用快速幂把乘法次数从 O($n$) 降到 O($\log n$)。
& 写快速幂递归和迭代版各一，验证 $3^{10}\bmod 7 = 4$，并说明二分拆指数的正确性。 \\[2mm]

*贡献计算
& 不枚举所有子结构，改为计算每个元素在最终答案中被计入的次数。常用于"所有子数组的 max/min/sum 之和""每个节点对多少条路径有贡献"等聚合题。核心公式：$\text{ans}=\sum_i \text{val}_i\times\text{count}_i$。
& 给定数组 $[2,1,3]$，用单调栈求每个元素作为最大值的左右边界，计算以它为最大值的子数组个数，验证所有子数组最大值之和是否为 $2\cdot1+1\cdot1+3\cdot3=12$。手算与代码一致。 \\[2mm]

组合数学基础
& 先判断分类是否互斥且完备，再使用加法原理；分步完成且步骤相互独立时使用乘法原理。确定计数对象与分类后，再判断是否需要 DP。
& 对"选一个人"与"先选组别再选组内成员"各写一种计数式，指出重复计数和漏计数的位置。 \\[2mm]

排列、组合与重复计数
& 区分排列 $P(n,k)$、组合 $C(n,k)$、可重复选择（$n^k$）和相同元素排列（多重集排列公式）。常与全排列、子集、频次分配结合。
& 手算 $C(10,3)$，分别解释"从 10 人中选 3 人排队"和"选 3 人开会"的答案为何不同。给 $\{a,a,b\}$ 写出所有不同排列。 \\[2mm]

组合递推与杨辉三角
& $C(n,k)=C(n-1,k-1)+C(n-1,k)$，边界为 $C(n,0)=C(n,n)=1$；大量询问时预处理、取模、滚动空间。
& 写杨辉三角前 $5$ 行，验证一行的和为 $2^n$，并说明组合数取模时不能随意使用实数除法（要用递推或乘法逆元）。 \\[2mm]

*多重计数、映射与双射
& 直接计数困难时，建立原问题与另一个易计数集合的双射（一一对应）；或按某个特征分类后对每类分别计数再加和；也可用"计数$\times$量=总量"的双视角技巧。
& 用隔板法（Stars and Bars）解释"长度为 $n$、和为 $S$ 的非负整数序列个数"为 $C(S+n-1,n-1)$，手算 $n=3,S=4$ 的结果并列举验证。举例说明如何把"子数组和为 $K$"映射为前缀差对 $(i,j)$ 满足 $\text{pre}[i]-\text{pre}[j]=K$。 \\[2mm]

一元幂方程与单调方程
& 求 $x^k=A$、$ax+b=c$ 或单调函数的整数根时，利用定义域、单调性和二分；一般高次多项式没有可直接套用的通用初等公式。
& 写整数二分求最大的 $x$ 使 $x^k\le A$，避免乘法溢出并回代检查；比较整数答案和实数近似的处理差异。 \\[2mm]

一元二次方程求解
& $ax^2+bx+c=0$（$a\neq 0$）：先算 $\Delta=b^2-4ac$。$\Delta<0$ 无实根；$\Delta\ge 0$ 继续。整数根条件：$\Delta$ 为完全平方数且 $-b\pm\sqrt{\Delta}$ 被 $2a$ 整除。非整数需精确输出：约分有理数，或提取 $\sqrt{\Delta}$ 的平方因子写成规范根式 $\frac{p\pm q\sqrt{s}}{r}$，保证 $\gcd(p,q,r)=1$ 且 $s$ 无平方因子。
& 写程序对 $x^2-3x+1=0$ 输出两个根的规范形式（含 $\sqrt{5}$），并处理 $a<0$、$b^2$ 溢出（用 \code{long long}）、$\Delta$ 非平方、分母约分和较大根选择。对照历年真题中 RSA 解密和一元二次方程输出题验证。 \\[2mm]

整数开方
& 需要精确整数答案或判完全平方数时，优先用整数运算并验证；不能把浮点 \code{sqrt} 的结果未经检查直接转整数。
& 写整数二分开方，避免 \code{mid*mid} 溢出，验证 $\lfloor\sqrt{10^{18}}\rfloor$。 \\[2mm]

二分与*三分
& 二分用于单调函数求零点或阈值：整数域上用闭区间不变量防死循环，实数域按精度或固定次数迭代。三分用于单峰/单谷函数求极值：每轮比较 $f(m_1)$ 与 $f(m_2)$ 扔掉极值不在的 $1/3$ 区间；离散域上区间缩小到 $\le 3$ 后枚举。与第 2 节"二分答案"的差异：这里面对的是具体函数 $f(x)$，不要求构造判定谓词。
& 写整数三分在 $[0,10]$ 上找 $f(x)=-(x-3)^2+5$ 的最大值，先用枚举验证极值位置，再解释三分为何每次扔掉 $1/3$ 区间，并说明整数三分何时需要转枚举收尾。 \\
\end{longtable}
\endgroup

\newpage
%==================================================================
\section{*数据范围与复杂度反推（OI 视角）}
%==================================================================

将总操作数展开为包含全部独立参数的公式，再与 $B=5\times10^8$ 比较。表中的 $\log$ 均按 $\log_2$；边界取满足公式不超过 $B$ 的最大整数，不计常数。$n=5000$ 时，O($n^2$) 为 $2.5\times10^7$，符合本时间口径；空间另行计算。

\begingroup
\small
\begin{longtable}{L{3.1cm} C{3.2cm} L{5.6cm} L{4.0cm}}
\toprule
复杂度 & 最大整数规模 & 边界校验 & 常见候选 \\
\midrule
\endfirsthead
\toprule
复杂度 & 最大整数规模 & 边界校验 & 常见候选 \\
\midrule
\endhead
\endfoot
\bottomrule
\endlastfoot
O($n^n$) & $n\le9$ & $9^9=387{,}420{,}489$，$10^{10}>B$ & 每个位置有 $n$ 种选择 \\[1.2mm]
O($n!$) & $n\le12$ & $12!=479{,}001{,}600$，$13!>B$ & 全排列 \\[1.2mm]
O($3^n$) & $n\le18$ & $3^{18}=387{,}420{,}489$ & 三进制状态 \\[1.2mm]
O($n^2 2^n$) & $n\le20$ & $20^2\cdot2^{20}=419{,}430{,}400$ & 状压双层转移 \\[1.2mm]
O($n2^n$) & $n\le24$ & $24\cdot2^{24}=402{,}653{,}184$ & 子集加末项 \\[1.2mm]
O($2^n$) & $n\le28$ & $2^{28}=268{,}435{,}456$，$2^{29}>B$ & 子集枚举 \\[1.2mm]
O($n^5$) & $n\le54$ & $54^5=459{,}165{,}024$ & 五重枚举 \\[1.2mm]
O($n^4$) & $n\le149$ & $149^4=492{,}884{,}401$ & 四重枚举 \\[1.2mm]
O($n^3$) & $n\le793$ & $793^3=498{,}677{,}257$ & 三重枚举 \\[1.2mm]
O($n^2\log n$) & $n\le6294$ & $6294^2\log_2 6294\approx4.999\times10^8$ & 二维状态加对数操作 \\[1.2mm]
O($n^2$) & $n\le22360$ & $22360^2=499{,}969{,}600$ & 点对、二维 DP \\[1.2mm]
O($n\sqrt n$) & $n\le629960$ & $n\sqrt n\approx4.99999\times10^8$ & 分块、因数枚举 \\[1.2mm]
O($n\log n$) & $n\le20{,}580{,}553$ & $n\log_2n\approx4.99999975\times10^8$ & 排序、堆、有序结构 \\[1.2mm]
O($n$) & $n\le500{,}000{,}000$ & $n=B$ & 扫描、前缀、线性 DP \\[1.2mm]
O($\log n$)/O(1) & 常用整数范围 & 仍须乘调用次数并计输入输出 & 二分定位、公式 \\
\end{longtable}
\endgroup

\textbf{多参数直接展开：}$T$ 组 O($n^2$) 写成 $Tn^2$；$Q$ 次长度 $n$ 的扫描写成 $Qn$；图写 $V+E$ 或 $(V+E)\log V$；背包写 $nW$；网格写 $rc$；状态图写 $V\cdot K$ 个状态。结果 $\le5\times10^8$ 即通过本时间口径，随后单独检查内存。


\newpage
%==================================================================
\section{*空间复杂度与四档内存反推}
%==================================================================

使用 $1\text{ MiB}=2^{20}\text{ B}$。下表只是一块连续载荷占满限制的数学上限；程序映像、其他数组、容器容量、分配器、临时对象和调用栈全部计入总内存，实装时必须逐项求和。按常见 64 位 GCC 环境估算：\code{char/bool}=1 B、\code{int}=4 B、\code{long long/double}=8 B；结构体用目标环境 \code{sizeof} 核准。

\begingroup
\small
\begin{longtable}{L{2.5cm} R{3.3cm} R{3.3cm} R{3.3cm} R{3.3cm}}
\toprule
限制 & 总字节/1 B 元素 & \code{int} 元素 & 8 B 元素 & 压位布尔 bit \\
\midrule
\endfirsthead
\toprule
限制 & 总字节/1 B 元素 & \code{int} 元素 & 8 B 元素 & 压位布尔 bit \\
\midrule
\endhead
\endfoot
\bottomrule
\endlastfoot
128 MiB & 134,217,728 & 33,554,432 & 16,777,216 & 1,073,741,824 \\[1.2mm]
256 MiB & 268,435,456 & 67,108,864 & 33,554,432 & 2,147,483,648 \\[1.2mm]
512 MiB & 536,870,912 & 134,217,728 & 67,108,864 & 4,294,967,296 \\[1.2mm]
1024 MiB & 1,073,741,824 & 268,435,456 & 134,217,728 & 8,589,934,592 \\
\end{longtable}

\begin{longtable}{L{3.8cm} C{3.1cm} C{3.1cm} C{3.1cm} C{3.1cm}}
\toprule
单一方阵理论边长 & 128 MiB & 256 MiB & 512 MiB & 1024 MiB \\
\midrule
\endfirsthead
\toprule
单一方阵理论边长 & 128 MiB & 256 MiB & 512 MiB & 1024 MiB \\
\midrule
\endhead
\endfoot
\bottomrule
\endlastfoot
\code{int a[n][n]} & $n\le5792$ & $n\le8192$ & $n\le11585$ & $n\le16384$ \\[1.2mm]
\code{long long a[n][n]} & $n\le4096$ & $n\le5792$ & $n\le8192$ & $n\le11585$ \\
\end{longtable}
\endgroup

\textbf{校准：}\code{int dp[5000][5000]} 是 100,000,000 B，约 95.37 MiB；两份约 190.73 MiB。$n=20000$ 的 O($n^2$) 时间是 $4\times10^8$，满足时间预算，但 \code{int[n][n]} 约 1525.88 MiB，1024 MiB 仍放不下。

\begin{longtable}{L{3.3cm} L{4.0cm} L{8.0cm}}
\toprule
结构/算法 & 渐进空间 & 必须展开的峰值 \\
\midrule
\endfirsthead
\toprule
结构/算法 & 渐进空间 & 必须展开的峰值 \\
\midrule
\endhead
\endfoot
\bottomrule
\endlastfoot
前缀和/差分 & O($n$) & 一个 long long 数组约 $8(n+1)$ B；原数组若仍存活要相加。 \\[1.5mm]
BFS/DFS & O($V+E$) & 图 + 标记/距离 + 队列/显式栈；递归另有 O(height) 调用栈。 \\[1.5mm]
邻接矩阵 & O($V^2$) & int 矩阵 $4V^2$ B。 \\[1.5mm]
数组邻接表 & O($V+E$) & 每弧三个 int 约 $12E+4V$ B；无向 $m$ 边存两倍弧，约 $24m+4V$ B。 \\[1.5mm]
Dijkstra & O($V+E$) & 图、距离和懒删除优先队列同时存活；旧状态也计峰值。 \\[1.5mm]
归并排序 & O($n$) & 原数组 + $sizeof(T)\cdot n$ 临时数组 + O($\log n$) 栈。 \\[1.5mm]
数组 Trie & O($L|\Sigma|$) & $4\cdot maxNode\cdot alphabet$ B；节点上界是总字符数加 1。 \\[1.5mm]
背包滚动 & O($W$) & 一个 int 层 $4(W+1)$ B；保留全部 $n$ 层则为 O($nW$)。 \\[1.5mm]
区间/二维 DP & O($n^2$) & long long 表为 $8n^2$ B；时间可行不表示表能开。 \\[1.5mm]
map/set/哈希 & O($n$) & 结点、指针、桶、对齐、负载因子和扩容峰值均超过裸值载荷。 \\
\end{longtable}

\begin{itemize}
  \item $\square$ 列出所有同时存活的大对象并求和，不仅计算最大的单个对象。
  \item $\square$ 分别在 128/256/512/1024 MiB 限制下判断方案，并指出首个超限对象。
  \item $\square$ 将无向边的两条弧、容器冗余、递归栈和临时数组计入峰值。
  \item $\square$ 判断 O($n^2$) 方案由时间还是空间先限制，并列出操作数和 MiB。
\end{itemize}

\newpage
%==================================================================
\section{*提交与编译检查}
%==================================================================

本节条目参考《常见错误文档表》中整理的实际代码错误。文件名、编译选项和判定结果以当年规则及题面为准。

\begingroup
\small
\begin{longtable}{L{2.7cm} L{5.4cm} L{6.4cm} L{1.8cm}}
\toprule
检查位置 & 常见错误 & 检查要求 & 可能结果 \\
\midrule
\endfirsthead
\toprule
检查位置 & 常见错误 & 检查要求 & 可能结果 \\
\midrule
\endhead
\endfoot
\bottomrule
\endlastfoot
文件读写 & 题目要求文件 I/O 时，文件名或扩展名错误；输入输出文件写反；重定向被注释；使用本机绝对路径 & 按题面逐字核对文件名；输入使用读取模式并重定向到标准输入，输出使用写入模式并重定向到标准输出；不得保留本机路径 & 无输出、读入失败、零分 \\[1.5mm]
入口与头文件 & \code{main}、\code{include} 或 \code{namespace} 拼写错误；头文件缺失；提交了错误语言格式 & 使用 C++14 编译最终提交文件，确认只有一个有效的 \code{main}，所用接口均有对应头文件 & CE \\[1.5mm]
字符串字面量 & 文件名等字符串误用单引号；字符串拼写与大小写不一致 & 字符使用单引号，字符串使用双引号；需要精确匹配的文本逐字符核对 & CE、读写失败 \\[1.5mm]
调试清理 & 删除调试输出时破坏括号、分号或条件结构；遗留额外输出 & 清理后重新格式化并完整编译；用样例核对输出中没有提示语、日志或多余空格行 & CE、WA \\[1.5mm]
粘贴与文件内容 & 重复粘贴函数或完整程序；把样例输入、大段表格或无关文本写入源码 & 搜索重复定义；检查提交文件类型和大小；源程序不得超过当年规定上限 & CE、提交失败 \\[1.5mm]
算法完整性 & 固定输出样例、按输入分支输出预置答案、随机输出或遗漏核心算法 & 使用未见过的合法小数据验证；确认输出由输入和算法状态计算得到 & WA、零分 \\[1.5mm]
平台接口 & 使用 \code{system()}、本机专用命令或依赖本地目录结构 & 仅使用比赛环境允许的标准接口；在规定编译环境运行最终版本 & CE、RE、无输出 \\[1.5mm]
初始化与多测 & 局部变量未初始化；数组或容器未在每组数据前复位 & 开启编译警告；逐项检查初值、不可达标记及多组数据的清空范围 & WA、UB \\
\end{longtable}
\endgroup

\begin{itemize}
  \item $\square$ 最终文件能以规定的 C++14 选项独立编译。
  \item $\square$ 文件读写仅在题面要求时启用，名称、扩展名和读写方向完全一致。
  \item $\square$ 已删除调试输出、固定答案、本机路径和平台专用调用。
  \item $\square$ 已核对提交文件、语言类型、源程序大小和多组数据初始化。
\end{itemize}

\newpage
%==================================================================
\section{*C++ UB 与实现风险}
%==================================================================

\begingroup
\small
\begin{longtable}{L{0.9cm} L{2cm} L{4.8cm} L{3.1cm} L{4.5cm}}
\toprule
种类 & 风险 & 错误示例 & 可能结果 & 修正 \\
\midrule
\endfirsthead
\toprule
种类 & 风险 & 错误示例 & 可能结果 & 修正 \\
\midrule
\endhead
\endfoot
\bottomrule
\endlastfoot

UB & 有符号溢出 & \code{int x=2000000000; x+=x;} & 结果不可依赖 & 使用 \code{long long} 或先判 \code{x>LIMIT-a} \\[2mm]
UB & 左移溢出 & \code{1<<31}（32 位 \code{int}） & 出人意料的值 & \code{1LL<<31} 或 \code{1u<<31} \\[2mm]
UB & 数组越界 & \code{int a[10]; a[10]=0;} & 覆盖相邻变量、段错误 & 每次访问前检查下标 \\[2mm]
UB & 整数除零 & \code{int x=a/b;}（\code{b} 可能为 $0$） & 运行时未定义 & 分母可能为零时先判 \\[2mm]
UB & 未初始化变量 & \code{int x; if(x)\{...\}} & 结果随机 & 定义时赋初值 \\[2mm]
UB & 迭代器失效 & \code{v.erase(it); it++;} & 跳过元素或崩溃 & \code{it=v.erase(it)} \\[2mm]
UB & 严格弱序不满足 & \code{sort(\dots,[](int a,int b){return a<=b;})} & 排序崩溃 & 改 \code{return a<b;} \\[2mm]
UB & 悬空引用 & \code{auto\& r=*v.end();} & 非法内存 & 不对 \code{end()} 解引用 \\[2mm]
资源 & 调用栈耗尽 & \code{void f(){f();}} & RE / SIGSEGV & 限制深度或改为迭代 \\[2mm]
UB & C++14 求值顺序 & \code{arr[i]=i++;} & 未定义行为（C++14） & 拆两行，先保存旧值 \\[2mm]
UB & int 乘后赋 long long & \code{long long x=a*b;} & 先在 int 中溢出 & \code{1LL*a*b} \\[2mm]
数值 & 浮点判等 & \code{double a,b; if(a==b)} & 相等判不出 & \code{fabs(a-b)<eps} 或改整数 \\[2mm]
性能 & 哈希退化 & \code{unordered\_map} 发生集中冲突 & 单次操作退化为 O($N$) & 仅将 O(1) 写作期望复杂度 \\[2mm]
UB & 基类析构非虚 & \code{Base*p=new Derived; delete p;} & 通过基类删除是未定义行为 & 基类析构函数声明为 \code{virtual} \\
\end{longtable}
\endgroup

{\small\noindent
\textbf{分类说明：}UB 不具有可依赖结果；未指定行为表示标准未规定具体选择；数值风险表示程序合法但数值可能错误；性能问题表示行为确定但复杂度可能超限。}

\newpage
%==================================================================
\section{*对拍流程}
%==================================================================

对拍由定义式暴力程序、候选程序、数据生成器和结果比较脚本组成。

\subsection*{三个独立程序}

\begin{enumerate}[itemsep=2pt]
  \item \textbf{brute.cpp}：按定义实现并优先保证正确；生成数据的规模须与其复杂度相匹配。
  \item \textbf{sol.cpp}：待提交的候选程序。
  \item \textbf{gen.cpp}：生成合法的小规模数据；以测试编号作为种子，保证失败用例可复现。
\end{enumerate}

\subsection*{分别编译}

任一源文件修改后均须重新编译：
{\small
\begin{verbatim}
g++ -std=c++14 -O2 -o brute brute.cpp
g++ -std=c++14 -O2 -o sol   sol.cpp
g++ -std=c++14 -O2 -o gen   gen.cpp
\end{verbatim}}

\subsection*{对拍脚本 run\_test.sh}

{\small
\begin{verbatim}
#!/bin/bash
for i in $(seq 1 1000); do
  ./gen "$i" > in.txt
  ./brute < in.txt > ans.txt
  ./sol   < in.txt > out.txt
  if ! diff -q ans.txt out.txt > /dev/null; then
    echo "WA on test $i"; cp in.txt hack.in; exit 1
  fi
done
echo "All passed"
\end{verbatim}}

流程：gen 生成输入 $\to$ brute 生成标准输出 $\to$ sol 生成候选输出 $\to$ diff 比较结果。出现差异时保存输入到 hack.in 并停止；对可能不终止的程序设置运行超时。

\subsection*{生成器模板 gen.cpp}

{\small
\begin{verbatim}
#include <cstdlib>
#include <iostream>
#include <random>
int main(int argc, char **argv) {
  unsigned seed = argc > 1
      ? static_cast<unsigned>(std::strtoul(argv[1], nullptr, 10))
      : 42u;                       // 无参数时仍可复现
  std::mt19937 rng(seed);
  std::uniform_int_distribution<int> n_dist(1, 10);
  std::uniform_int_distribution<int> a_dist(0, 99);
  int n = n_dist(rng);
  std::cout << n << "\n";
  for (int i = 0; i < n; i++)
    std::cout << a_dist(rng)
              << (i+1 == n ? '\n' : ' ');
  return 0;
}
\end{verbatim}}

模板仅演示可复现的随机种子。实际测试还须按题意加入极值、重复值、结构化数据和无解数据，均匀随机数据不能替代这些用例。

\subsection*{边界用例}

\begin{itemize}
  \item 最小规模（$n=1$）和最大规模（但仅能暴力跑的范围）
  \item 全相等 / 全 0 / 全 1
  \item 已升序 / 已降序
  \item 无解输入（$-1$ 或 NO 的情况）
  \item 紧贴边界（$n=10^9$ 触顶、模数边界等）
\end{itemize}

\subsection*{检查项}

\begin{itemize}
  \item $\square$ 能在限定时间内完成暴力程序、候选程序、生成器和比较脚本。
  \item $\square$ 生成器使用固定种子，并将首个差异输入保存到 \code{hack.in}。
  \item $\square$ 出现 WA 时，能够分别核对生成器、暴力程序、候选程序和比较脚本。
\end{itemize}

\newpage
%==================================================================
\section{*重点补充项}
%==================================================================

\begin{itemize}
  \item 历年 CSP-J/NOIP 普及组真题多次考查模拟和扫描，同时改变语言边界、整数范围与输入格式。
  \item \textbf{需要重点核对的内容}：
        \begin{enumerate}
          \item 字符串区间哈希：能够预处理前缀哈希，并按统一区间定义比较子串。
          \item 数学与计数：能够使用贡献计算、映射或双射；区分排列与组合、加法原理与乘法原理；完成判别式与根式规范化。
          \item 复杂度反推：在实现前将全部参数代入时间预算，并独立核算峰值内存。
          \item 对拍：准备定义式暴力、可复现生成器、边界用例和首个失败用例。
        \end{enumerate}
  \item 线性筛、二分三分、Trie、区间哈希、树状数组、单调队列、状态图未必以裸模板出现，但常作为综合题的一部分。
\end{itemize}

\medskip
\section*{创作声明与许可}
\addcontentsline{toc}{section}{创作声明与许可}

除另有说明的第三方内容外，本文由 \texttt{scmmm} 原创。本文以
\href{https://creativecommons.org/licenses/by-nc/4.0/deed.zh-hans}{Creative Commons 署名--非商业性使用 4.0 国际版（CC BY-NC 4.0）}
发布。

在注明原作者为 \texttt{scmmm}、保留本许可声明和许可证链接，并标明修改内容的前提下，允许非商业性转载、分享与修改。任何商业性质的使用、转载、修改、发行，或将本文用于商业产品、课程或服务，均须事先获得 \texttt{scmmm} 的授权。

\end{document}
